% !PRINTED 351
\weberpage{351}

Um die imaginären Wurzeln zu erhalten, ersetzt man $\mathrm{tg}\,\varphi$ durch $\varrho\,\mathrm{tg}\,\varphi$, wenn $\varrho$ eine imaginäre dritte Einheitswurzel bedeutet.

b) Im zweiten Falle bestimmen wir den Winkel $\vartheta$, gleichfalls im ersten Quadranten, aus

\weq{7}{\frac{g}{2} = \sqrt{\frac{e^3}{27}}\,\frac{1}{\sin\vartheta}}

\[
\sqrt{R} = \sqrt{\frac{e^3}{27}}\,\mathrm{cotg}\,\vartheta
\]

\[
x = \sqrt{\frac{e}{3}}\left(\sqrt[3]{\mathrm{cotg}\,\frac{\vartheta}{2}} + \sqrt[3]{\mathrm{tg}\,\frac{\vartheta}{2}}\right),
\]

\weq{8}{\mathrm{tg}\,\frac{\vartheta}{2} = \mathrm{tg}\,\varphi^3,}

\weq{9}{x = 2\sqrt{\frac{e}{3}}\,\frac{1}{\sin 2\varphi}.}

c) Im letzten Falle endlich, wo drei reelle Wurzeln vorhanden sind, setzen wir

\weq{10}{\frac{g}{2} = \sqrt{\frac{e^3}{27}}\cos\vartheta}

\[
\sqrt{R} = i\sqrt{\frac{e^3}{27}}\sin\vartheta
\]

\[
x = \sqrt{\frac{e}{3}}\left(\sqrt[3]{\cos\vartheta + i\sin\vartheta} + \sqrt[3]{\cos\vartheta - i\sin\vartheta}\right),
\]

oder nach dem \gesperrt{Moivre'schen} Satze:

\weq{11}{x = 2\sqrt{\frac{e}{3}}\cos\frac{1}{3}\vartheta.}

Nimmt man $\vartheta$ wieder im ersten Quadranten, so erhält man für die beiden anderen gleichfalls reellen Wurzeln

\[
-2\sqrt{\frac{e}{3}}\cos\frac{\pi + \vartheta}{3}, \qquad -2\sqrt{\frac{e}{3}}\cos\frac{\pi - \vartheta}{3}.
\]

Alle diese Formeln sind für die logarithmische Rechnung eingerichtet.
